\documentclass[
    ref = refDemo,
    style = 0,
    sectionpagestyle = 1
]{spBeamer}
\spTitle{\texttt{spBeamer} Demo}
\spAuthor[SP]{Sweet Pastry}
\spAffiliation[FDU]{Fudan University, Shanghai, China}

\begin{document}
    \section{Introduction}
        \begin{frame}{Logo}
            % \begin{figure}
                % \centering
                \includegraphics[width=\textwidth]{cover.pdf}
                % \caption{Cover, Sweet Pastry}
            % \end{figure}
        \end{frame}

    \section{Math}
        \subsection{subsection1}
            \begin{frame}{\texttt{uncover} command}
                \uncover<+->{%
                \begin{definition}[Linear Functional]
                    A linear functional on a vector space \( X \) over \( \mathbb{R} \) or \( \mathbb{C} \) is a map \( f : X \to \mathbb{R} \) (or \( \mathbb{C} \)) such that:
                    \[
                    f(\alpha x + \beta y) = \alpha f(x) + \beta f(y) \quad \forall x,y \in X, \ \alpha, \beta \in \mathbb{R} \ (\text{or } \mathbb{C})
                    \]
                \end{definition}
                \uncover<+->{
                \begin{definition}[Sublinear Functional]
                    A map \( p: X \to \mathbb{R} \) is called sublinear if:
                    \[
                    p(x + y) \leq p(x) + p(y) \quad \text{and} \quad p(\lambda x) = \lambda p(x) \quad \forall x,y \in X, \ \lambda \geq 0
                    \]
                \end{definition}
                }}
            \end{frame}

            \begin{frame}
                \begin{theorem}[Hahn-Banach, Normed Space Version]
                    Let \( X \) be a normed space, \( M \) a subspace of \( X \), and \( f: M \to \mathbb{R} \) a bounded linear functional. Suppose \( p: X \to \mathbb{R} \) is a sublinear functional such that:
                    \[
                    f(x) \leq p(x) \quad \forall x \in M
                    \]
                    Then, there exists an extension \( F: X \to \mathbb{R} \) of \( f \) such that:
                    \[
                    F(x) \leq p(x) \quad \forall x \in X
                    \]
                    and \( F \) is linear and bounded.
                \end{theorem}
            \end{frame}

        \subsection{subsection2}
            \begin{frame}
                \uncover<+->{%
                The proof uses Zorn's Lemma and proceeds by extending \( f \) step by step to larger subspaces. The key idea is to define:
                \uncover<+->{%
                \[
                F(x + \alpha y) = F(x) + \alpha f(y)
                \]
                \uncover<+->{%
                for some \( y \notin M \) and ensure the extension satisfies the sublinear constraint:
                \uncover<+->{%
                \[
                F(x + y) \leq p(x + y)
                \]
                \uncover<+->{%
                Applying Zorn's Lemma to the collection of all extensions leads to the existence of the desired functional \( F \).
                }}}}}
            \end{frame}

    \section{chem}
        \subsection{Chem1}
            \begin{frame}
                \begin{center}
                    \begin{tikzpicture}
                        \node at (0.78,2.95) {\chemfig{
                                <[:-25,1.176]
                                (
                                    -[0,0.01,,,line width=6.1pt,shorten <=-1.5pt,shorten >=-1.5pt]
                                    >[:25,1.176]
                                    -[:135,0.8] (
                                            -[:30,1]-[:-35,1]OH
                                            )
                                    -[180,1]
                                    -[:-135,0.8]
                                )}};
                        \node at (0,0.3){\chemfig{
                                    -[:25,1.176]
                                (
                                    -[:-25,1.176]
                                    <[:-135,0.8]
                                    -[180,,,,line width=4.8pt,shorten <=-2.2pt,shorten >=-2.2pt]
                                    >[:135,0.8]
                                )}};
                        \draw [dashed]  (0,1.75) -- (0,2.9);
                        \draw [dashed]  (0,0.3) -- (0,1.3);
                        \draw (0,0.3) ellipse (0.75cm and 0.35cm);
                        \draw (0,2.85) ellipse (0.75cm and 0.35cm);
                        \node at (0,1.5) {\Large Fe};
                    \end{tikzpicture}
                \end{center}
            \end{frame}
\end{document}